% ==============================================================================
%  1.2  Relación entre variables. Relaciones lineales y linealizables.
%       Asociación y causalidad.
% ==============================================================================
\section{Relación entre variables}
\label{sec:t1-relacion}

\subsection{Relación funcional y relación estadística}

\begin{definicion}[Relación funcional y relación estadística]
  \begin{itemize}
    \item Una \textbf{relación funcional}, expresada mediante una fórmula matemática, es una
      relación perfecta entre variables: el valor de las variables explicativas determina
      exactamente el valor de la respuesta.
    \item Una \textbf{relación estadística} es algo intermedio entre una relación funcional y
      la ausencia de relación: conocidas las variables explicativas, la respuesta no puede
      predecirse exactamente, pero sí de forma aproximada.
  \end{itemize}
\end{definicion}

El interés está en modelizar relaciones estadísticas, mediante modelos del tipo
\begin{equation}\label{eq:t1-modelo-general}
  Y = m(X_1,\ldots,X_k) + \varepsilon ,
\end{equation}
donde $m(X_1,\ldots,X_k)$ es el \emph{patrón sistemático} de la relación, en torno al cual se
distribuyen los datos con cierta variabilidad que recoge el \emph{error aleatorio}
$\varepsilon$. A este error se le atribuyen una serie de propiedades (distribución, media,
varianza\ldots).

\begin{figure}[H]
  \centering
  \begin{subfigure}[t]{0.48\linewidth}
    \centering
    \begin{tikzpicture}
      \begin{axis}[width=6cm, height=5cm, xlabel={Unidades vendidas ($X$)},
          ylabel={Ventas en \$ ($Y$)}, xmin=0, xmax=160, ymin=0, ymax=320]
        \addplot[recta, domain=0:150] {2*x};
        \addplot[puntos, mark size=1.5pt] coordinates {(25,50) (75,150) (130,260)};
        \node[font=\footnotesize, anchor=west] at (axis cs:95,120) {$Y=2X$};
      \end{axis}
    \end{tikzpicture}
    \caption{Relación funcional: los puntos están exactamente sobre la recta.}
  \end{subfigure}\hfill
  \begin{subfigure}[t]{0.48\linewidth}
    \centering
    \begin{tikzpicture}
      \begin{axis}[width=6cm, height=5cm, xlabel={Edad ($X$)},
          ylabel={Nivel de esteroides ($Y$)}, xmin=7, xmax=26, ymin=0, ymax=30]
        \addplot[recta, domain=8:25, samples=50] {-26.3+4.84*x-0.1165*x^2};
        \addplot[puntos, mark size=1.3pt] coordinates {
          (8,2) (8,6) (9,9.5) (9,12) (9,4) (10,11) (11,15) (12,11) (12,19) (13,13)
          (13,21) (14,21) (14,24) (16,22) (17,19) (17,24) (17,25) (18,21) (18,20.5)
          (19,22.5) (21,20) (21,24) (23,23) (23,27) (24.5,26) (25,22) (25,20)};
      \end{axis}
    \end{tikzpicture}
    \caption{Relación estadística: los puntos se distribuyen alrededor de un patrón
      sistemático (en este caso, curvo).}
  \end{subfigure}
  \caption{Relación funcional frente a relación estadística.}
\end{figure}

\begin{definicion}[Variables relacionadas][def:variables-relacionadas]
  Dos variables $X$ e $Y$ están \textbf{relacionadas} si la distribución condicional de
  $Y \cond X=x$ cambia al cambiar $x$. En caso contrario, las variables son
  \textbf{independientes}.
\end{definicion}

Uno de los objetivos de la modelización estadística es describir de forma sencilla cómo
cambian las distribuciones condicionales de $Y$ al cambiar $X$. La misma idea se extiende a la
relación entre $Y$ y un conjunto de variables explicativas $X_1,\ldots,X_k$.

\subsection{Selección de las variables explicativas}

Una parte esencial del ajuste de un modelo es la \textbf{selección de las variables
explicativas} $X_1,\ldots,X_k$, especialmente cuando $k$ es muy grande en relación con el
tamaño muestral $n$. Los criterios de selección son:
\begin{itemize}
  \item \textbf{Criterio principal}: en qué medida contribuyen a explicar la variabilidad de la
    respuesta. Son preferibles las que explican más.
  \item \textbf{Conocimiento del fenómeno}: por ejemplo, su importancia como agente causal.
  \item \textbf{Cuestiones prácticas de la toma de datos}: la precisión con la que puede
    observarse la variable y la facilidad, rapidez o coste de cada observación.
\end{itemize}

\subsection{Forma de la relación: relaciones lineales y linealizables}

Otra decisión importante es la \textbf{forma de la relación}. Normalmente no se conoce a priori
y debe decidirse empíricamente a partir de los datos. Las funciones lineales suelen ser una
buena elección por varias razones:
\begin{itemize}
  \item su \emph{simplicidad}, que es la más evidente, aunque no la única;
  \item muchas \emph{relaciones no lineales} pueden \textbf{linealizarse} mediante
    transformaciones de $X$, de $Y$ o de ambas (inversa, logaritmo, cuadrado, etc.);
  \item las relaciones de forma muy compleja pueden aproximarse mediante \textbf{funciones
    lineales ``a trozos''}.
\end{itemize}
Todo ello amplía el campo de aplicación de los modelos lineales y los convierte en una
herramienta fundamental para el ajuste de modelos estadísticos.

\begin{figure}[H]
  \centering
  \begin{tikzpicture}
    \begin{axis}[width=6.5cm, height=4.2cm, xmin=0, xmax=1.05, ymin=0, ymax=1.1,
        xtick={0,0.2,0.4,0.6,0.8,1},
        xticklabels={$x_0{=}0$,$x_1$,$x_2$,$x_3$,$x_4$,$x_5{=}1$},
        ytick=\empty, axis lines=left, tick label style={font=\scriptsize}]
      \addplot[NavyBlue, thick, domain=0:1, samples=60] {4*x*(1-x)};
      \addplot[red, thick, mark=none] coordinates
        {(0,0) (0.2,0.64) (0.4,0.96) (0.6,0.96) (0.8,0.64) (1,0)};
      \foreach \x in {0.2,0.4,0.6,0.8} {
        \edef\tmp{\noexpand\draw[dashed, gray] (axis cs:\x,0) -- (axis cs:\x,{4*\x*(1-\x)});}
        \tmp }
    \end{axis}
  \end{tikzpicture}
  \caption{Aproximación de una relación no lineal (azul) mediante una función lineal a
    trozos (rojo).}
\end{figure}

\subsection{Asociación y causalidad}
\label{sec:t1-causalidad}

Una relación puede ser \emph{débil}, en cuyo caso puede carecer de importancia práctica. Pero,
\textbf{aunque exista una relación estadística fuerte entre $X$ e $Y$, no puede asumirse que
$X$ sea la causa de $Y$} (relación causa--efecto):
\begin{itemize}
  \item la relación observada puede deberse a la influencia de terceras variables
    (\textbf{variables confusoras}), cuyos efectos solo pueden neutralizarse mediante el
    \emph{diseño de experimentos} (\cref{sec:t1-datos});
  \item puede existir una relación causal en sentido contrario, de $Y$ hacia $X$;
  \item hay incluso situaciones en las que la relación no es interpretable.
\end{itemize}

\begin{ejemplo}[Chocolate y premios Nobel][ej:t1-chocolate]
  En 2012 F.~H.~Messerli publicó en el \emph{New England Journal of Medicine} el estudio
  ``El consumo de chocolate, la función cognitiva y los Premios Nobel''. Con datos por
  países, encontró una relación muy fuerte ($r=0.791$, $p<0.0001$) entre el consumo anual de
  chocolate per cápita y el número de premios Nobel por cada 10 millones de habitantes.

  Poco después aparecieron otros trabajos (Maurage, Heeren y Pesenti, 2013, \emph{The Journal
  of Nutrition}) que señalaban otros índices relacionados con el número de premiados:
  \begin{center}
    \small
    \begin{tabular}{lc}
      \toprule
      Relación entre países & $r$ \\
      \midrule
      Nobel -- consumo de té            & $0.03$ \\
      Nobel -- consumo de vino          & $0.16$ \\
      Nobel -- tiendas IKEA por 10 M hab. & $0.82$ \\
      Nobel -- PIB per cápita           & $0.73$ \\
      Consumo de chocolate -- PIB per cápita & $0.66$ \\
      \bottomrule
    \end{tabular}
  \end{center}
  El número de tiendas de IKEA está más relacionado con el número de premios Nobel que el
  consumo de chocolate.

  Esta relación no permite concluir que comer chocolate aumente la probabilidad de obtener un
  premio Nobel. La tabla apunta a una \textbf{variable confusora}, la riqueza del país (PIB per
  cápita), relacionada tanto con el consumo de chocolate ($r=0.66$) como con el número de
  premios Nobel ($r=0.73$). Además, se trata de datos observacionales agregados por países,
  con los que no pueden establecerse relaciones causa--efecto (\cref{sec:t1-datos}).
\end{ejemplo}
